\documentclass[10pt,a4paper]{amsart}
\usepackage[margin=19mm]{geometry}
\usepackage{amsmath,amssymb,mathtools}
\usepackage[scr=rsfs,cal=euler]{mathalfa}
\usepackage{xfrac}
\usepackage[unicode,hidelinks,bookmarks=false]{hyperref}
\newcommand{\ie}{i.e.\ }
\newcommand{\cf}{cf.\ }
\newcommand{\resp}{respectively\ }
\newcommand{\Subsection}{\subsection}
\providecommand{\nobreakdash}{\nobreak\hbox{-}}
\setlength{\emergencystretch}{2em}
\multlinegap0pt
\usepackage{lmodern}
\usepackage{mathtools,bm}
\usepackage{booktabs}
\usepackage{array}
\usepackage{float}
\usepackage{enumitem}
\usepackage{xurl}
\usepackage{cleveref}
\newtheorem{theorem}{Theorem}[section]
\newtheorem{lemma}[theorem]{Lemma}
\newtheorem{proposition}[theorem]{Proposition}
\newtheorem{corollary}[theorem]{Corollary}
\newtheorem{definition}[theorem]{Definition}
\newtheorem{remark}[theorem]{Remark}
\newcommand{\F}{\mathbb{F}}
\newcommand{\C}{\mathcal{C}}
\newcommand{\Kcal}{\mathcal{K}}
\newcommand{\Lcal}{\mathcal{L}}
\newcommand{\Dcal}{\mathcal{D}}
\newcommand{\supp}{\operatorname{supp}}
\newcommand{\wt}{\operatorname{wt}}
\newcommand{\poly}{\operatorname{poly}}
\newcommand{\SCL}{\operatorname{SCL}}
\newcommand{\TV}{\operatorname{TV}}
\newcommand{\E}{\mathbb{E}}
\newcommand{\Prb}{\mathbb{P}}
\newcommand{\ket}[1]{\lvert #1\rangle}
\newcommand{\bra}[1]{\langle #1\rvert}
\newcommand{\abs}[1]{\left\lvert #1\right\rvert}
\newcommand{\norm}[1]{\left\lVert #1\right\rVert}
\newcommand{\one}{\mathbf{1}}
\newcommand{\RS}{\operatorname{RS}}
\newcommand{\GRS}{\operatorname{GRS}}
\newcommand{\QFT}{\operatorname{QFT}}
\newcommand{\Had}{\operatorname{Had}}
\newcommand{\id}{\mathrm{id}}
\begin{document}
\title[Efficient Exact Quantum Sampling from the Sun-Wootters Distr]{Efficient Exact Quantum Sampling from the Sun-Wootters Distribution for Optimal Polynomial Intersection}
\author{Sunghyeon Jo}
\date{}
\maketitle
\noindent\emph{Source and licence.} Efficient Exact Quantum Sampling from the Sun-Wootters Distribution for Optimal Polynomial Intersection. Version arXiv:2607.16541v2 (CC BY 4.0). This material is distributed under the Creative Commons Attribution 4.0 International licence, http://creativecommons.org/licenses/by/4.0/, as recorded in the arXiv OAI licence metadata for this exact version and shown on the abstract page https://arxiv.org/abs/2607.16541v2. Full rights notice: licenses/arxiv-2607.16541-CC-BY-4.0.txt.\par\smallskip
\noindent\textit{Editorial scope.} Counted unit: the original \texttt{theorem} environment \texttt{thm:SW-import} at source lines 281--317 together with its complete proof at lines 319--321; the delivered document keeps source lines 175--776 so that every referenced label and every citation resolves inside it. Native editable source is the paper's own concatenated TeX; the AMS wrapper is editorial formatting only. No publisher class, upstream script or unrelated artwork is redistributed.
\par\smallskip\noindent\textit{Reference note.} This article ships no compiled bibliography (biblatex); the entries delivered here are composed from the author's own BibTeX (.bib) fields, with no field invented and no entry dropped.
\section{Preliminaries}

\subsection{Fourier Analysis over a Prime Field}

Let $p$ be prime and let $\omega=e^{2\pi i/p}$.  For $f:\F_p\to\mathbb C$, define
\begin{equation}
   \widehat f(z):=\frac1p\sum_{a\in\F_p} f(a)\omega^{-az},
   \qquad
   f(a)=\sum_{z\in\F_p}\widehat f(z)\omega^{az}.
   \label{eq:fourier-convention}
\end{equation}
Parseval's identity is
\begin{equation}
   \sum_{z\in\F_p}\abs{\widehat f(z)}^2
   =\frac1p\sum_{a\in\F_p}\abs{f(a)}^2.
   \label{eq:parseval-one-dimensional}
\end{equation}
For $s,x\in\F_p^n$, write $\chi_s(x)=\omega^{\langle s,x\rangle}$.  Expectations over a finite set are uniform unless a distribution is displayed explicitly.

\subsection{Balanced OPI}

Let $\alpha_1,\ldots,\alpha_m\in\F_p$ be distinct, with $m\le p$, and let
\begin{equation}
   B=\begin{pmatrix}
      1&\alpha_1&\cdots&\alpha_1^{n-1}\\
      \vdots&\vdots&&\vdots\\
      1&\alpha_m&\cdots&\alpha_m^{n-1}
   \end{pmatrix}\in\F_p^{m\times n}.
   \label{eq:vandermonde}
\end{equation}
For $x\in\F_p^n$, the vector $Bx$ consists of evaluations of the polynomial with coefficient vector $x$.  Suppose that all input sets have a common density
\[
   \abs{S_i}=\rho p,
   \qquad 0<\rho<1.
\]
Define the satisfaction ratio
\begin{equation}
   s(x):=\frac1m\sum_{i=1}^m\one\{\langle b_i,x\rangle\in S_i\},
   \label{eq:score}
\end{equation}
where $b_i$ is row $i$ of $B$, and define the normalized discrepancy function
\begin{equation}
   g_i(a):=\frac{\one\{a\in S_i\}-\rho}{\sqrt{\rho(1-\rho)}}.
   \label{eq:gi}
\end{equation}
Then
\begin{equation}
   \frac1p\sum_a g_i(a)=0,
   \qquad
   \frac1p\sum_a g_i(a)^2=1,
   \qquad
   \widehat g_i(0)=0,
   \qquad
   \sum_{z\ne0}\abs{\widehat g_i(z)}^2=1.
   \label{eq:gi-properties}
\end{equation}
For $0\le k\le m$, set
\begin{equation}
   q_k(x):=\sum_{\substack{T\subseteq[m]\\\abs{T}=k}}
              \prod_{i\in T}g_i(\langle b_i,x\rangle).
   \label{eq:qk}
\end{equation}

We study asymptotic families with $m\to\infty$, $1\le n<m\le p$, $n/m\to r=2\mu\in(0,1]$, prime $p=p(m)$, and $\rho=\rho_m=1/2+o(1)$.  All asymptotic constants may depend on the fixed limiting parameters but not on the particular input sets.  For probability distributions $P,Q$ on the same finite set, we use
\[
   \TV(P,Q):=\frac12\sum_x\abs{P(x)-Q(x)}.
\]

\paragraph{Computational model.}
The integers $p,m,n$, the distinct evaluation points, and the common cardinality $\rho p$ are classical input data.  The sets are accessed by the coherent membership oracle in \cref{eq:membership-oracle}.  A deterministic classical preprocessor may use the public input data and the requested accuracy to compute field constants, the right inverse $R$, the multipliers $\nu_i$, and a description of the quantum circuit; its polynomial running time is included in the total complexity.  The resulting oracle circuit is uniform in the standard sense that this description is produced in time polynomial in the classical input length.

Put $b=\lceil\log_2p\rceil$ and identify $\F_p$ with the code subspace $\mathcal H_p=\operatorname{span}\{\ket0,\ldots,\ket{p-1}\}\subseteq(\mathbb C^2)^{\otimes b}$.  Every reversible field operation is specified as a permutation on all $2^b$ binary strings; encodings $p,\ldots,2^b-1$ have fixed total behavior and are never populated by the ideal preparation.  Likewise, an approximate prime-modulus QFT is a qubit unitary whose restriction to $\mathcal H_p$ is within the stated operator-norm error of the mathematical $p$-dimensional QFT, with an arbitrary unitary extension on $\mathcal H_p^\perp$.  Running time counts one- and two-qubit gates, reversible bit operations, and oracle calls.  Arbitrary rotations and the prime-modulus QFT are synthesized over a fixed universal gate set to the precision specified in \cref{cor:finite-precision}.

\subsection{The Sun--Wootters Target Distribution}

Fix $\delta\in[0,1/2-\mu]$ and define
\begin{equation}
   \ell=\lfloor(\mu+\delta)m\rfloor,
   \qquad
   \sigma=\lfloor\log\log\ell\rfloor,
   \qquad
   \Kcal=\{\ell-\sigma,\ldots,\ell\}\cap\{0,\ldots,m\}.
   \label{eq:shell}
\end{equation}
Choose
\begin{equation}
   u_k:=\binom{m}{k}^{-1/2}\one\{k\in\Kcal\}.
   \label{eq:weights}
\end{equation}
The unnormalized target amplitude and its probability distribution are
\begin{equation}
   F(x):=\sum_{k\in\Kcal}u_kq_k(x),
   \qquad
   P_u(x):=\frac{\abs{F(x)}^2}{\sum_{z\in\F_p^n}\abs{F(z)}^2}.
   \label{eq:Pu}
\end{equation}
Because every $g_i$ is real, $F(x)$ is real, so this agrees with the square used by Sun and Wootters.

For $0\le t\le1/2$, the balanced semicircle function is
\begin{equation}
   \SCL_{1/2}(t)=\frac12+\sqrt{t(1-t)}.
   \label{eq:scl}
\end{equation}

We use the following consequence of the proof of Sun and Wootters's Theorem~1.6, specifically their denominator expansion, Lemma~4.11, and Lemma~5.2 \cite{SunWootters2026}.


\begin{theorem}[Sun--Wootters analytic estimate]
\label{thm:SW-import}
Let $r=2\mu$, let $\rho_m=1/2+o(1)$, and use \cref{eq:shell,eq:weights,eq:Pu}.  Define the binary entropy with natural logarithms by
\[
   h(t)=-t\log t-(1-t)\log(1-t),
\]
and
\begin{align}
   E(\mu,\delta)
   &:=(\mu+\delta)h\!\left(\frac{\mu}{\mu+\delta}\right)
      +(1-\mu-\delta)h\!\left(\frac{\mu}{1-\mu-\delta}\right)
      -h(2\mu),
      \label{eq:E-exponent}\\
   G(\mu,\lambda)
   &:=(1-2\mu)\log2+h(2\mu)
      -(4\mu-1)h\!\left(\frac{\lambda}{4\mu-1}\right)\notag\\
   &\hspace{2.8em}
      -(2-4\mu)h\!\left(\frac{2\mu-\lambda}{2-4\mu}\right)
      +\lambda\log\!\left(\frac2\pi\right),
      \label{eq:G-exponent}
\end{align}
where parameters are restricted to the domain on which the entropy arguments lie in $[0,1]$.
If there is a $\lambda$ such that
\begin{equation}
   E(\mu,\delta)+G(\mu,\lambda)<0,
   \label{eq:SW-condition}
\end{equation}
then there exist constants $c>0$ and $C<\infty$ such that, uniformly over all balanced input sets,
\begin{align}
   \E_{x\in\F_p^n}\abs{F(x)}^2
      &=\abs{\Kcal}+O(m^Ce^{-cm}),
      \label{eq:SW-denominator}\\
   \E_{x\sim P_u}s(x)
      &\ge \SCL_{1/2}(\mu+\delta)-o(1).
      \label{eq:SW-score}
\end{align}
\end{theorem}

\begin{proof}[Derivation from the cited estimates]
Sun and Wootters define the same discrepancy functions and the same distribution in their equations~(4.2), (4.17), and (4.18); the proportionality constant in their $u_k$ is immaterial, and we fix it as in \cref{eq:weights}.  Their Proposition~4.7 evaluates the weight-zero part of the denominator as exactly $\sigma+1=\abs{\Kcal}$.  Under \cref{eq:SW-condition}, their Lemma~5.2, combined with Lemma~4.10 and the proof of Theorem~1.6, bounds every nonzero-weight correction by an exponentially decaying term.  Summing the polynomially many weight and shell-index choices gives the additive $O(m^Ce^{-cm})$ term in \cref{eq:SW-denominator}; this is the denominator displayed in their equation~(4.35).  Their Lemma~4.11 then gives the expectation bound under the distribution $P_u$.  The strict exponent margin absorbs the $o(m)$ terms caused by $\rho_m=1/2+o(1)$ and by integer rounding.  All Fourier estimates in that argument are worst-case bounds over the sets $S_i$, so the constants are uniform over the promised inputs.
\end{proof}


\begin{lemma}[Stability under sublinear parameter perturbations]
\label{lem:SW-stability}
The conclusions of \cref{thm:SW-import} remain valid if $n/(2m)=\mu+o(1)$ and the largest active weight satisfies $\ell/m=\mu+\delta+o(1)$, provided the limiting parameters obey the strict inequality \cref{eq:SW-condition}.  The same is true for any fixed $\lambda$ in the interior of its entropy domain.
\end{lemma}

\begin{proof}
In the derivation above, every binomial coefficient contributes its entropy exponent plus $o(m)$, the balanced Fourier estimate contributes $G(\mu,\lambda)+o(1)$ per coordinate, and the shell width is $o(m)$.  Replacing $\mu$ and $\ell/m-\mu$ by convergent sequences changes the continuous entropy exponents by $o(1)$.  A strict negative limiting margin therefore remains negative by a fixed constant for all sufficiently large $m$, and the same summations yield \cref{eq:SW-denominator,eq:SW-score}.
\end{proof}

\begin{remark}
Theorem~1.6 of Sun and Wootters is stated as an existence theorem.  The stronger distributional conclusion in \cref{eq:SW-score} is what their proof establishes before taking a maximizing point: their Lemma~4.11 directly lower-bounds expectation under $P_u$.  Equation \eqref{eq:SW-denominator} is the denominator estimate appearing in their equation~(4.35), after the common normalization of the $u_k$ is fixed as in \cref{eq:weights}.
\end{remark}

\section{The Fourier-Domain Target State}

Let
\begin{equation}
   H:=B^{\mathsf T}\in\F_p^{n\times m},
   \qquad \C:=\operatorname{im}(B),
   \qquad
   \C^\perp:=\ker H.
   \label{eq:H-code}
\end{equation}
For $y\in\F_p^m$, define
\begin{equation}
   a_y:=u_{\wt(y)}\prod_{i:y_i\ne0}\widehat g_i(y_i),
   \label{eq:ay}
\end{equation}
where $u_k=0$ outside $\Kcal$.  Define
\begin{equation}
   Z_0:=\sum_y\abs{a_y}^2,
   \qquad
   \ket{\Psi_0}:=\frac1{\sqrt{Z_0}}\sum_y a_y\ket{y}\ket{Hy}.
   \label{eq:labeled-state}
\end{equation}
For each syndrome $s\in\F_p^n$, let
\begin{equation}
   \Lcal_s:=\{y\in\F_p^m:Hy=s,\ \wt(y)\in\Kcal\},
   \qquad
   A_s:=\sum_{y\in\Lcal_s}a_y,
   \label{eq:fibers}
\end{equation}
and set
\begin{equation}
   Z_I:=\sum_s\abs{A_s}^2,
   \qquad
   \ket{\widehat I}:=\frac1{\sqrt{Z_I}}\sum_sA_s\ket{s}.
   \label{eq:ideal-syndrome}
\end{equation}

\begin{lemma}[Fourier expansion]
\label{lem:fourier-expansion}
For every $x\in\F_p^n$,
\begin{equation}
   F(x)=\sum_{y\in\F_p^m}a_y\chi_{Hy}(x)
       =\sum_{s\in\F_p^n}A_s\chi_s(x).
   \label{eq:fourier-expansion}
\end{equation}
Consequently,
\begin{equation}
   Z_I=\E_{x\in\F_p^n}\abs{F(x)}^2,
   \label{eq:ZI-parseval}
\end{equation}
and applying the inverse QFT over $\F_p^n$ to \cref{eq:ideal-syndrome} produces a state whose measurement distribution is exactly $P_u$.
\end{lemma}

\begin{proof}
By Fourier inversion and $\widehat g_i(0)=0$,
\begin{align*}
   q_k(x)
   &=\sum_{\abs{T}=k}\prod_{i\in T}
       \left(\sum_{z_i\ne0}\widehat g_i(z_i)
       \omega^{z_i\langle b_i,x\rangle}\right)\\
   &=\sum_{\wt(y)=k}
       \left(\prod_{i:y_i\ne0}\widehat g_i(y_i)\right)
       \omega^{\langle B^{\mathsf T}y,x\rangle}.
\end{align*}
Multiplying by $u_k$ and summing over $k$ proves the first equality in \cref{eq:fourier-expansion}; grouping terms by $s=Hy$ proves the second.  Character orthogonality gives
\[
   \frac1{p^n}\sum_x\abs{F(x)}^2
   =\sum_s\abs{A_s}^2=Z_I.
\]
Finally, under the convention
\[
   (\QFT_p^{-1})^{\otimes n}\ket{s}
      =p^{-n/2}\sum_x\chi_s(x)\ket{x},
\]
the amplitude at $x$ is $F(x)/(p^{n/2}\sqrt{Z_I})$.  After normalization, its squared magnitude gives the probability in \cref{eq:Pu}.
\end{proof}

\begin{lemma}[Norm of the labeled error state]
\label{lem:Z0}
For the weights in \cref{eq:weights},
\begin{equation}
   Z_0=\abs{\Kcal}=\sigma+1
\end{equation}
for all sufficiently large $m$.
\end{lemma}

\begin{proof}
Fix a support $T\subseteq[m]$.  By \cref{eq:gi-properties},
\[
   \sum_{\supp(y)=T}
      \prod_{i\in T}\abs{\widehat g_i(y_i)}^2
   =\prod_{i\in T}\sum_{z\ne0}\abs{\widehat g_i(z)}^2=1.
\]
For weight $k$, there are $\binom mk$ supports and $u_k^2=\binom mk^{-1}$, so every active weight contributes one.
\end{proof}

\section{Complete List Decoding of Syndrome Fibers}

\subsection{The Dual Code Is Generalized Reed--Solomon}

For distinct $\alpha_1,\ldots,\alpha_m$, define
\begin{equation}
   \nu_i:=\left(\prod_{j\ne i}(\alpha_i-\alpha_j)\right)^{-1}.
   \label{eq:dual-multipliers}
\end{equation}

\begin{lemma}[Dual evaluation code]
\label{lem:dual-GRS}
The code $\C^\perp=\ker B^{\mathsf T}$ is
\begin{equation}
   \C^\perp
   =\left\{(\nu_1f(\alpha_1),\ldots,\nu_mf(\alpha_m)):
       f\in\F_p[X],\ \deg f<m-n\right\}.
   \label{eq:dual-GRS}
\end{equation}
In particular it is a generalized Reed--Solomon code of dimension $m-n$ and distance $n+1$.
\end{lemma}

\begin{proof}
Let $q\in\F_p[X]$ have degree at most $m-2$.  Lagrange interpolation writes
\[
   q(X)=\sum_{i=1}^m q(\alpha_i)
        \prod_{j\ne i}\frac{X-\alpha_j}{\alpha_i-\alpha_j}.
\]
The coefficient of $X^{m-1}$ on the left is zero, while on the right it is
$\sum_i\nu_iq(\alpha_i)$.  Therefore
\begin{equation}
   \sum_{i=1}^m\nu_iq(\alpha_i)=0
   \quad\text{whenever }\deg q\le m-2.
   \label{eq:lagrange-orthogonality}
\end{equation}
If $h$ has degree less than $n$ and $f$ has degree less than $m-n$, then $q=hf$ has degree at most $m-2$.  Equation \eqref{eq:lagrange-orthogonality} shows that the vector in \cref{eq:dual-GRS} is orthogonal to every evaluation vector $(h(\alpha_i))_i$, so it lies in $\C^\perp$.  The displayed family has dimension $m-n$, equal to $\dim\C^\perp$, hence equality holds.  The distance statement is the standard root bound: a nonzero polynomial of degree less than $m-n$ has at most $m-n-1$ zeros among the evaluation points.
\end{proof}

\subsection{Fibers as Decoding Lists}

Since $B$ has full column rank, $H=B^{\mathsf T}$ has full row rank.  Choose any right inverse $R\in\F_p^{m\times n}$ satisfying $HR=I_n$, and define the canonical representative
\begin{equation}
   r_s:=Rs.
   \label{eq:representative}
\end{equation}
Then
\begin{equation}
   Hy=s
   \quad\Longleftrightarrow\quad
   y=r_s-c\text{ for a unique }c\in\C^\perp.
   \label{eq:coset-bijection}
\end{equation}
Consequently,
\begin{equation}
   \Lcal_s
   =\{r_s-c:c\in\C^\perp,\ \wt(r_s-c)\in\Kcal\}.
   \label{eq:fiber-list}
\end{equation}
Thus a complete decoder for $\C^\perp$ centered at $r_s$ through radius $\ell$, followed by shell filtering, returns exactly the fiber $\Lcal_s$.

Polynomial-time list decoding up to the Johnson radius originates with Guruswami and Sudan \cite{GuruswamiSudan1999}.  For a uniform reversible implementation with complexity polynomial in $\log p$, we use the following deterministic theorem.

\begin{theorem}[Chatterjee--Harsha--Kumar \cite{ChatterjeeEtAl2026}]
\label{thm:CHK}
For every finite field $\F$, block length $M$, and dimension $K<M$, there is a deterministic algorithm running in $\poly(M,\log\abs{\F})$ bit operations that outputs every Reed--Solomon codeword agreeing with a received word in more than $\sqrt{(K-1)M}$ coordinates.
\end{theorem}

\begin{corollary}[Complete decoding of the dual OPI code]
\label{cor:dual-list-decoding}
There is a deterministic algorithm that, on input $s\in\F_p^n$, outputs exactly $\Lcal_s$ in time $\poly(m,\log p)$ whenever
\begin{equation}
   m-\ell>\sqrt{(m-n-1)m}.
   \label{eq:finite-Johnson}
\end{equation}
The output list has polynomial size.
\end{corollary}

\begin{proof}
By \cref{lem:dual-GRS}, divide coordinate $i$ of $r_s$ by the nonzero multiplier $\nu_i$.  This maps $\C^\perp$ to the ordinary Reed--Solomon code of dimension $m-n$ and preserves Hamming distance.  Under \cref{eq:finite-Johnson}, every codeword at distance at most $\ell$ has agreement at least $m-\ell>\sqrt{(m-n-1)m}$, so \cref{thm:CHK} outputs all of them.  Rescale, form $y=r_s-c$, retain exactly those with $\wt(y)\in\Kcal$, remove duplicates, and sort lexicographically.  Every step is deterministic and polynomial time.  A polynomial-time machine can write only polynomially many output entries, which gives the list-size bound.
\end{proof}

\begin{lemma}[Explicit uniform list bound]
\label{lem:explicit-list-bound}
Let $\C$ be an $[m,k]$ generalized Reed--Solomon code, put $a=m-\ell$, and suppose $a^2>m(k-1)$.  Every Hamming ball of radius $\ell$ contains at most
\begin{equation}
   J(m,k,a):=
   \frac{a-(k-1)}{a^2/m-(k-1)}\le m^2
   \label{eq:explicit-list-bound}
\end{equation}
codewords.  If $k/m$ and $a/m$ converge while $a^2/m^2-k/m$ stays positive by a fixed constant, then $J(m,k,a)=O(1)$.
\end{lemma}

\begin{proof}
Since $a\le m$ and $a^2>m(k-1)$, we have $a>k-1$; in particular, the numerator and denominator in \cref{eq:explicit-list-bound} are positive.
Suppose that $N$ codewords lie in the ball.  For each codeword choose an $a$-element subset $S_u$ of coordinates on which it agrees with the received word, and put $t_i=\abs{\{u:i\in S_u\}}$.  Distinct generalized Reed--Solomon codewords agree in at most $k-1$ coordinates, so
\[
   \sum_i\binom{t_i}{2}
   =\sum_{u<v}\abs{S_u\cap S_v}
   \le\binom N2(k-1).
\]
Since $\sum_it_i=Na$, Cauchy--Schwarz gives
\[
   \sum_i\binom{t_i}{2}
   \ge\frac12\left(\frac{N^2a^2}{m}-Na\right).
\]
Rearranging yields
$N(a^2/m-(k-1))\le a-(k-1)$.  The displayed bound also holds when $N=1$ because
$J-1=a(1-a/m)/(a^2/m-(k-1))\ge0$.  Finally, $a^2-m(k-1)$ is a positive integer, so the denominator in \cref{eq:explicit-list-bound} is at least $1/m$, while its numerator is at most $m$.  Under a fixed normalized margin, the numerator and denominator are both linear in $m$, proving the last assertion.
\end{proof}

\section{Coherent Fiber Summation}

\subsection{A Canonical Fixed-Length List Representation}

Fix one deterministic implementation of \cref{cor:dual-list-decoding}, including lexicographic sorting and duplicate removal.  Apply \cref{lem:explicit-list-bound} with $k=m-n$ and $a=m-\ell$, and set
\begin{equation}
   B_\ell:=\left\lceil J(m,m-n,m-\ell)\right\rceil\le m^2,
   \qquad
   L:=2^{\lceil\log_2B_\ell\rceil}<2m^2,
   \qquad q:=\log_2L.
   \label{eq:global-L}
\end{equation}
Represent the sorted fiber as
\[
   \Lcal_s=(y_{s,0},\ldots,y_{s,N_s-1})
\]
and pad slots $N_s,\ldots,L-1$ with a validity bit equal to zero.  The same $L$ is used for every syndrome.  At a fixed strict asymptotic Johnson margin, both $B_\ell$ and $L$ are constant in $m$.

\begin{lemma}[Canonical reversible indexing]
\label{lem:index-unitary}
There is a uniform quantum circuit whose size is polynomial in $m$ and $\log p$.  Its action on every supported basis state is
\begin{equation}
   \ket{y_{s,j}}\ket{s}\ket{0^q}_J\ket0_Q\ket{0^{\mathrm{anc}}}
   \longmapsto
   \ket{\mathbf 0}_{\F_p^m}\ket{s}\ket{j}_J\ket1_Q\ket{0^{\mathrm{anc}}},
   \qquad 0\le j<N_s.
   \label{eq:index-map}
\end{equation}
Here $Q$ is a match flag.  The circuit is a unitary on the full computational basis, including inputs that are not valid fiber elements or valid field encodings.
\end{lemma}

\begin{proof}
Pad the deterministic decoder and its sorting routine to a fixed polynomial number of steps.  Bennett's compute--copy--uncompute construction turns this fixed-time computation into a reversible circuit with polynomial overhead \cite{Bennett1973}.  On input $s$, compute the padded canonical list and its validity bits into workspace $W$.  Reversibly compare $y$ with every valid slot, compute the first matching index into $J$, and set $Q=1$ exactly when a match exists.  When no slot matches, use the harmless convention $J=0,Q=0$.  Uncompute all comparison scratch space before changing the error register.  Thus a genuine match in slot zero is distinguished from a no-match input.

Controlled on $Q=1$, use $(W,J)$ to compute the selected list entry into a temporary register, subtract it from the error register, and reverse the selection computation to clear the temporary register.  Finally reverse the decoder computation to clear $W$; its input $s$ has not changed.  If $y=y_{s,j}$ is valid, the error register is zero and the retained pair is $(J,Q)=(j,1)$.  On a no-match or invalid-encoding input the controlled subtraction is skipped.  Every step is a reversible permutation of basis states with specified behavior on all invalid inputs, so the circuit is a unitary on the entire Hilbert space and \cref{eq:index-map} holds on the support of \cref{eq:labeled-state}.
\end{proof}

\subsection{Uniform List-Index Projection}

Apply \cref{lem:index-unitary} to \cref{eq:labeled-state}.  The result is
\begin{equation}
   \frac1{\sqrt{Z_0}}
   \sum_s\sum_{j<N_s}a_{y_{s,j}}\ket{\mathbf 0}_{\F_p^m}\ket{s}\ket{j}_J\ket1_Q.
   \label{eq:indexed-state}
\end{equation}
Apply $\Had^{\otimes q}$ to $J$ and measure $(J,Q)$ in the computational basis.

\begin{theorem}[Exact coherent fiber sum]
\label{thm:exact-fiber-sum}
If $Z_I>0$, then conditioned on the joint outcome $J=0^q,Q=1$, the syndrome register in \cref{eq:indexed-state} is exactly $\ket{\widehat I}$ from \cref{eq:ideal-syndrome}.  The success probability is exactly
\begin{equation}
   p_{\mathrm{fib}}=\frac{Z_I}{LZ_0}.
   \label{eq:success-exact}
\end{equation}
After the inverse QFT over $\F_p^n$, measurement samples exactly from $P_u$.
\end{theorem}

\begin{proof}
For every $j\in\{0,\ldots,L-1\}$,
\[
   \langle0^q\vert \Had^{\otimes q}\vert j\rangle=L^{-1/2}.
\]
Hence the unnormalized successful branch is
\begin{align*}
   \frac1{\sqrt{LZ_0}}\sum_s\sum_{j<N_s}a_{y_{s,j}}\ket{s}
   &=\frac1{\sqrt{LZ_0}}\sum_sA_s\ket{s}.
\end{align*}
Its squared norm is $Z_I/(LZ_0)$, and normalization gives \cref{eq:ideal-syndrome}.  This is a valid probability: by Cauchy--Schwarz in each fiber,
\[
   Z_I=\sum_s\abs{\sum_{y\in\Lcal_s}a_y}^2
      \le\sum_sN_s\sum_{y\in\Lcal_s}\abs{a_y}^2
      \le LZ_0.
\]
The final claim follows from \cref{lem:fourier-expansion}.
\end{proof}

\begin{remark}[Uniform padding across lists]
A syndrome-dependent projection length $N_s$ would produce the state
$\sum_s A_sN_s^{-1/2}\ket{s}$.  Padding every list to one global length $L$ replaces these varying factors by the common factor $L^{-1/2}$ and yields the Sun--Wootters state.
\end{remark}

\subsection{Algorithm Statement}

One sampling attempt consists of the following operations.
\begin{enumerate}[label=\textbf{\arabic*.},leftmargin=2.2em]
  \item Prepare the labeled state $\ket{\Psi_0}$ in \cref{eq:labeled-state}.
  \item Apply the reversible canonical indexing circuit of \cref{lem:index-unitary}.
  \item Apply $\Had^{\otimes q}$ to the list-index register and measure $(J,Q)$.  Restart unless $(J,Q)=(0^q,1)$.
  \item Apply $(\QFT_p^{-1})^{\otimes n}$ to the syndrome register and measure $x\in\F_p^n$.
\end{enumerate}

The next section proves that all four steps have polynomial complexity in a standard random-access model and that the restart probability is polynomially bounded under the Sun--Wootters condition.

\section{State Preparation, Complexity, and Robustness}

\subsection{Preparing the Labeled Error State}

We first isolate the same local data-access primitive used by DQI.

\begin{definition}[Fourier-state access]
For each $i\in[m]$, let $U_i$ be a clean unitary satisfying
\begin{equation}
   U_i\ket{0}=\sum_{z\in\F_p}\widehat g_i(z)\ket{z}.
   \label{eq:local-fourier-state}
\end{equation}
Controlled applications of $U_i$ and $U_i^\dagger$ are available at cost $T_g$.
\end{definition}

\begin{proposition}[Labeled error-state preparation]
\label{prop:shell-prep}
Given Fourier-state access, the labeled state $\ket{\Psi_0}$ can be prepared with
$\poly(m,\log p)+O(mT_g)$ elementary operations and clean ancillas.
\end{proposition}

\begin{proof}
Prepare a uniform superposition over $k\in\Kcal$ using a standard reversible range-state construction, with all rotations synthesized to the precision required below. Conditioned on $k$, prepare the Dicke state
\[
   \binom mk^{-1/2}\sum_{\substack{t\in\{0,1\}^m\\\wt(t)=k}}\ket{t}.
\]
A uniform Dicke state can be prepared sequentially: if $N$ positions and $R$ ones remain, rotate the next bit with amplitudes $\sqrt{(N-R)/N}$ and $\sqrt{R/N}$ and update the counters.  Induction on $N$ shows that every weight-$k$ string receives amplitude $\binom mk^{-1/2}$; reversible fixed-point arithmetic gives polynomial gate complexity.

For each coordinate $i$, controlled on $t_i=1$, apply $U_i$ to a field register initially equal to zero.  Because $\widehat g_i(0)=0$, the resulting field value is nonzero exactly on the branch $t_i=1$.  Compute the predicate $\one\{y_i\ne0\}$ and xor it into $t_i$, clearing the whole support register.  Compute $\wt(y)$ reversibly, subtract it from the $k$ register, and then reverse the weight computation; this clears both the $k$ register and its counter workspace.  The amplitude of a resulting vector $y$ of weight $k$ is
\[
   \frac1{\sqrt{\abs{\Kcal}}}\binom mk^{-1/2}
   \prod_{i:y_i\ne0}\widehat g_i(y_i)
   =\frac{a_y}{\sqrt{Z_0}},
\]
where \cref{lem:Z0} was used.  Finally compute $Hy$ by reversible field arithmetic.  For finite precision, expand the range-state, Dicke-state, controlled local-state, and prime-QFT routines into a total of $G_{\mathrm{app}}=\poly(m,\log p,\log(1/\varepsilon))$ approximate primitive unitaries.  Synthesize each primitive to operator error at most $\varepsilon/G_{\mathrm{app}}$; the hybrid argument then bounds the total prepared-state error by $\varepsilon$.  Exact reversible arithmetic does not consume this error budget.
\end{proof}

\subsection{From Membership Access to Fourier-State Access}

Suppose the input is supplied through the standard random-access membership oracle
\begin{equation}
   O_S:\ket{i}\ket{a}\ket{b}
   \longmapsto
   \ket{i}\ket{a}\ket{b\oplus\one\{a\in S_i\}}.
   \label{eq:membership-oracle}
\end{equation}

\begin{proposition}[Balanced local-state preparation]
\label{prop:membership-to-fourier}
Assume $\rho\in[\rho_*,1-\rho_*]$ for a constant $\rho_*>0$ and that $\rho$ is known.  Given $O_S$ and $O_S^\dagger$, a controlled version of \cref{eq:local-fourier-state} can be implemented to operator-norm error at most $\varepsilon$ using $O_{\rho_*}(1)$ oracle calls and $\poly(\log p,\log(1/\varepsilon))$ additional gates.
\end{proposition}

\begin{proof}
Because every coordinate has the same known density $\rho$, the phase choices below are independent of $i$; the entire construction can therefore be controlled coherently by an index register.  Exact amplitude amplification with a known marked fraction prepares the uniform states
\[
   \ket{S_i}=\frac1{\sqrt{\rho p}}\sum_{a\in S_i}\ket{a},
   \qquad
   \ket{\overline S_i}=\frac1{\sqrt{(1-\rho)p}}\sum_{a\notin S_i}\ket{a}
\]
using $O(1/\sqrt\rho)$ and $O(1/\sqrt{1-\rho})$ membership queries, respectively \cite{BrassardEtAl2002}.  Prepare a branch qubit as
$\sqrt{1-\rho}\ket0-\sqrt\rho\ket1$ and, conditioned on the branch, prepare $\ket{S_i}$ or $\ket{\overline S_i}$.  The joint state is
\[
   \sum_{a\in S_i}\frac{\sqrt{1-\rho}}{\sqrt{\rho p}}\ket0\ket a
   -\sum_{a\notin S_i}\frac{\sqrt\rho}{\sqrt{(1-\rho)p}}\ket1\ket a.
\]
The branch is $0$ precisely on $S_i$ and $1$ precisely on its complement.  Using the branch as the oracle target maps both cases to $\ket1$; one Pauli $X$ then clears it to $\ket0$.  The remaining field register is
\[
   \frac1{\sqrt p}\sum_a g_i(a)\ket a.
\]
Applying the negative-phase QFT over $\mathbb Z_p$ gives \cref{eq:local-fourier-state} by \cref{eq:fourier-convention}.  Exact amplitude amplification uses efficiently specifiable phase rotations; approximating those rotations and the arbitrary-modulus QFT to total operator error $\varepsilon$ costs $\poly(\log p,\log(1/\varepsilon))$ gates \cite{HalesHallgren2000}.  The common density $\rho$ and the chosen phase schedule determine the branch phases.  A controlled one-qubit phase correction preserves the relative minus sign before the branches are combined.  The query count is constant because $\rho$ is bounded away from zero and one.
\end{proof}

\subsection{Polynomial Success Probability}

\begin{lemma}[Postselection probability]
\label{lem:norm-transfer}
Under the assumptions of \cref{thm:SW-import},
\begin{equation}
   \frac{Z_I}{Z_0}=1+O\!\left(\frac{m^Ce^{-cm}}{\sigma+1}\right)=1+o(1).
   \label{eq:norm-ratio}
\end{equation}
Consequently, for all sufficiently large $m$,
\begin{equation}
   p_{\mathrm{fib}}\ge\frac1{2L}.
   \label{eq:success-lower-bound}
\end{equation}
\end{lemma}

\begin{proof}
By \cref{lem:fourier-expansion}, $Z_I=\E_x\abs{F(x)}^2$.  Apply \cref{eq:SW-denominator} and \cref{lem:Z0}, then substitute into \cref{eq:success-exact}.
\end{proof}

Thus the required numbers of attempts in the ideal and finite-precision circuits are, respectively,
\[
   N_{\mathrm{ideal}}=\left\lceil2L\log(1/\beta)\right\rceil,
   \qquad
   N_{\mathrm{app}}=\left\lceil4L\log(1/\beta)\right\rceil.
\]
They produce a successful list-index projection with probability at least $1-\beta$, giving a standard bounded-error algorithm with restart.

\subsection{Finite-Precision Robustness}

\begin{lemma}[Conditioning stability]
\label{lem:conditioning-stability}
Let $\ket\psi$ and $\ket\phi$ be unit vectors with $\norm{\ket\psi-\ket\phi}\le\varepsilon$, and let $\Pi$ be a projector.  Write $p=\norm{\Pi\ket\psi}^2$ and $q=\norm{\Pi\ket\phi}^2$.  If $\varepsilon<\sqrt p$, then
\begin{align}
   \abs{p-q}&\le2\varepsilon,
   \label{eq:prob-stability}\\
   \norm{\frac{\Pi\ket\psi}{\sqrt p}-
          \frac{\Pi\ket\phi}{\sqrt q}}
      &\le\frac{2\varepsilon}{\sqrt p}.
   \label{eq:conditional-stability}
\end{align}
\end{lemma}

\begin{proof}
The reverse triangle inequality gives
$\abs{\sqrt p-\sqrt q}\le\norm{\Pi(\ket\psi-\ket\phi)}\le\varepsilon$.
Since $\sqrt p,\sqrt q\le1$, this implies \cref{eq:prob-stability}.  For the normalized vectors,
\begin{align*}
  \norm{\frac{\Pi\ket\psi}{\sqrt p}-
         \frac{\Pi\ket\phi}{\sqrt q}}
  &\le
    \frac{\norm{\Pi(\ket\psi-\ket\phi)}}{\sqrt p}
    +\sqrt q\abs{\frac1{\sqrt p}-\frac1{\sqrt q}}\\
  &\le\frac{\varepsilon}{\sqrt p}
       +\frac{\abs{\sqrt q-\sqrt p}}{\sqrt p}
   \le\frac{2\varepsilon}{\sqrt p}.
\end{align*}
\end{proof}

\begin{corollary}[Finite-precision implementation]
\label{cor:finite-precision}
Under \cref{eq:success-lower-bound}, an implementation whose premeasurement state is within
\begin{equation}
   \varepsilon\le
   \min\left\{\frac1{8L},\frac{\eta}{4\sqrt{2L}}\right\}
   \label{eq:precision-choice}
\end{equation}
of the ideal state immediately before the fiber-projection measurement, and whose final inverse QFT has operator error at most $\eta/2$, has success probability at least $1/(4L)$ and, conditioned on success, produces an output distribution at total variation distance at most $\eta$ from $P_u$.
\end{corollary}


\begin{thebibliography}{99}
\bibitem[]{Bennett1973} Charles H. Bennett. Logical Reversibility of Computation. IBM Journal of Research and Development 17 (1973) 525--532
\bibitem[]{BrassardEtAl2002} Gilles Brassard; Peter H{\o}yer; Michele Mosca; Alain Tapp. Quantum Amplitude Amplification and Estimation. Quantum Computation and Information 305 (2002) 53--74
\bibitem[]{ChatterjeeEtAl2026} Soham Chatterjee; Prahladh Harsha; Mrinal Kumar. Deterministic List Decoding of {Reed--Solomon} Codes. Proceedings of the 58th Annual ACM Symposium on Theory of Computing (2026) 2040--2051
\bibitem[]{GuruswamiSudan1999} Venkatesan Guruswami; Madhu Sudan. Improved Decoding of {Reed--Solomon} and Algebraic-Geometric Codes. IEEE Transactions on Information Theory 45 (1999) 1757--1767
\bibitem[]{HalesHallgren2000} Lisa Hales; Sean Hallgren. An Improved {Quantum Fourier Transform} Algorithm and Applications. Proceedings of the 41st Annual Symposium on Foundations of Computer Science (2000) 515--525
\bibitem[]{SunWootters2026} Yihang Sun; Mary Wootters. On Worst-Case Optimal Polynomial Intersection. (2026)
\end{thebibliography}
\end{document}
